\documentclass[10pt, letterpaper]{article}

% Packages:
\usepackage[
    ignoreheadfoot,
    top=0.4 cm,
    bottom=0.4 cm,
    left=0.8 cm,
    right=0.8 cm,
    footskip=0.2 cm,
]{geometry}
\usepackage{titlesec}
\usepackage{tabularx}
\usepackage{array}
\usepackage[dvipsnames]{xcolor}
\definecolor{primaryColor}{RGB}{0, 0, 0}
\usepackage{enumitem}
\usepackage{amsmath}
\usepackage[
    pdftitle={Tanmay Bhuskute's CV},
    pdfauthor={Tanmay Bhuskute},
    pdfcreator={LaTeX with RenderCV},
    colorlinks=true,
    urlcolor=primaryColor
]{hyperref}
\usepackage[pscoord]{eso-pic}
\usepackage{calc}
\usepackage{bookmark}
\usepackage{lastpage}
\usepackage{changepage}
\usepackage{paracol}
\usepackage{ifthen}
\usepackage{needspace}
\usepackage{iftex}

\ifPDFTeX
    \input{glyphtounicode}
    \pdfgentounicode=1
    \usepackage[T1]{fontenc}
    \usepackage[utf8]{inputenc}
    \usepackage{lmodern}
\fi

\usepackage{XCharter}

% Some settings:
\raggedright
\AtBeginEnvironment{adjustwidth}{\partopsep0pt}
\pagestyle{empty}
\setcounter{secnumdepth}{0}
\setlength{\parindent}{0pt}
\setlength{\topskip}{0pt}
\setlength{\columnsep}{0.15cm}
\pagenumbering{gobble}

% ── Section title ──
\titleformat{\section}{\needspace{2\baselineskip}\bfseries\large}{}{0pt}{}[\vspace{1pt}\titlerule]
\titlespacing{\section}{-1pt}{0.15 cm}{0.08 cm}

\renewcommand\labelitemi{$\vcenter{\hbox{\small$\bullet$}}$}

% ── Bullet lists ──
\newenvironment{highlights}{
    \begin{itemize}[
        topsep=0.10 cm,
        parsep=0.06 cm,
        partopsep=0pt,
        itemsep=0.06 cm,
        leftmargin=12pt
    ]
}{
    \end{itemize}
}

\newenvironment{highlightsforbulletentries}{
    \begin{itemize}[
        topsep=0.10 cm,
        parsep=0.06 cm,
        partopsep=0pt,
        itemsep=0.04 cm,
        leftmargin=12pt
    ]
}{
    \end{itemize}
}

\newenvironment{onecolentry}{
    \begin{adjustwidth}{0 cm + 0.00001 cm}{0 cm + 0.00001 cm}
}{
    \end{adjustwidth}
}

\newenvironment{twocolentry}[2][]{
    \onecolentry
    \def\secondColumn{#2}
    \setcolumnwidth{\fill, 4.5 cm}
    \begin{paracol}{2}
}{
    \switchcolumn \raggedleft \secondColumn
    \end{paracol}
    \endonecolentry
}

\newenvironment{threecolentry}[3][]{
    \onecolentry
    \def\thirdColumn{#3}
    \setcolumnwidth{, \fill, 4.5 cm}
    \begin{paracol}{3}
    {\raggedright #2} \switchcolumn
}{
    \switchcolumn \raggedleft \thirdColumn
    \end{paracol}
    \endonecolentry
}

\newenvironment{header}{
    \setlength{\topsep}{0pt}\par\kern\topsep\centering\linespread{0.6}
}{
    \par\kern\topsep
}

\let\hrefWithoutArrow\href

\begin{document}
    \newcommand{\AND}{\unskip
        \cleaders\copy\ANDbox\hskip\wd\ANDbox
        \ignorespaces
    }
    \newsavebox\ANDbox
    \sbox\ANDbox{$|$}

    \begin{header}
        \fontsize{16 pt}{24 pt}\selectfont Tanmay Bhuskute

        \vspace{3 pt}

        \normalsize
        \mbox{ +1 (732) 858-3027 }%
        \kern 1.0 pt%
        \AND%
        \kern 1.0 pt%
        \mbox{ tanmay.v.bhuskute@gmail.com }%
        \kern 1.0 pt%
        \AND%
        \kern 1.0 pt%
        \mbox{ github.com/Tanmayb05 }%
        \kern 1.0 pt%
        \AND%
        \kern 1.0 pt%
        \mbox{ linkedin.com/in/tanmay-bhuskute }%
        \kern 1.0 pt%
        \AND%
        \kern 1.0 pt%
        \mbox{ Austin, TX }%
    \end{header}

    \vspace{2 pt}

%═══════════════════════════════════════
\section{SUMMARY}
%═══════════════════════════════════════
\begin{onecolentry}
Infrastructure and Software Engineer with 2 years operating large-scale cloud distributed systems on AWS at Siemens. Built Kubernetes platforms, CI/CD automation, and Spark data pipelines with Java, Python, Helm, and Terraform; strengths in toil elimination, cost optimization, observability, and applying ML/GenAI to operational efficiency.
\end{onecolentry}

%═══════════════════════════════════════
\section{EDUCATION}
%═══════════════════════════════════════
\begin{twocolentry}{
    Aug 2024 - Dec 2026
}
    \textbf{Master of Science in Computer Science} \\
    University of Cincinnati
\end{twocolentry}
\begin{onecolentry}
Relevant Coursework: Advanced Algorithms, Intelligent Data Analysis, Database Theory.
\end{onecolentry}
\vspace{0.15 cm}

\begin{twocolentry}{
    Aug 2018 - Jul 2022
}
    \textbf{Bachelor of Engineering in Computer Engineering} \\
    AISSMS College of Engineering
\end{twocolentry}
\begin{onecolentry}
Relevant Coursework: Data Mining, Data Analytics, Object-Oriented Programming.
\end{onecolentry}
%═══════════════════════════════════════
\section{TECHNICAL SKILLS}
%═══════════════════════════════════════
\begin{onecolentry}
\textbf{Programming:} Java, Python, Go, C/C++, SQL, JavaScript, TypeScript, Bash, functional and object-oriented design.
\end{onecolentry}
\vspace{0.06 cm}
\begin{onecolentry}
\textbf{Data Systems:} Apache Spark, PySpark, Delta Lake, medallion pipelines, streaming \& batch data processing, ETL, REST APIs.
\end{onecolentry}
\vspace{0.06 cm}
\begin{onecolentry}
\textbf{Infrastructure \& Cloud:} AWS (EC2, S3, EFS, Lambda, SQS), multi-cloud (Azure, GCP), Kubernetes, ArgoCD, Helm, CRDs, Terraform, Ansible, Vault.
\end{onecolentry}
\vspace{0.06 cm}
\begin{onecolentry}
\textbf{Reliability \& Observability:} GitLab CI/CD, Jenkins, Docker, Linux, CloudWatch, Datadog, SonarQube, Blue-Green deploys.
\end{onecolentry}
\vspace{0.06 cm}
\begin{onecolentry}
\textbf{ML \& GenAI for Ops:} Gemini, Vertex AI, LangChain, PyTorch, scikit-learn, pandas, NumPy.
\end{onecolentry}

%═══════════════════════════════════════
\section{WORK EXPERIENCE}
%═══════════════════════════════════════
\begin{twocolentry}{Jul 2022 -- Jun 2024}
\textbf{Software Development Engineer} \,|\, Siemens Industry Software
\end{twocolentry}
\begin{highlights}
  \item Operated \textbf{40+ large-scale production environments} on AWS, executing \textbf{44 upgrades} in 4 months and troubleshooting live distributed-system issues across EC2, Linux, databases, and containerized services.
  \item Built internal automation frameworks in Python (Boto3) that replaced legacy Ruby scripts, eliminating deployment toil and cutting deployment time by \textbf{40\%} across AWS environments.
  \item Migrated \textbf{40+ environments} from process-based deployments to \textbf{Kubernetes} using ArgoCD, Helm, CRDs, GitLab CI/CD, Amazon EFS, and HashiCorp Vault, driving adoption of scalable deployment patterns.
  \item Designed \textbf{Ansible Blue-Green deployment} framework with automated rollback, reducing production upgrade downtime by \textbf{50\%} and enabling full recovery within 15 minutes.
  \item Cut monthly AWS spend \textbf{40\%} (\$15K to \$9K) via EC2 rightsizing and automated cleanup of orphaned resources, using CloudWatch and Datadog observability to drive cost and reliability decisions; also raised automated code coverage from \textbf{35\% to 78\%} as SonarQube Specialist.
\end{highlights}

\vspace{0.15 cm}

\begin{twocolentry}{Feb 2026 -- Present}
\textbf{Graduate Researcher} \,|\, University of Cincinnati
\end{twocolentry}
\begin{highlights}
  \item Designed two-stage memory-augmented \textbf{LLM pipeline} applying GenAI (Gemini, Vertex AI) to cut operational latency \textbf{25.7x} (39.4s vs 1012.1s per run) while preserving task fidelity across 53 runs on 20+ apps.
\end{highlights}

\vspace{0.15 cm}

\begin{twocolentry}{Mar 2022 -- Jun 2022}
\textbf{Software Developer Intern} \,|\, Siemens Industry Software
\end{twocolentry}
\begin{highlights}
  \item Automated AWS provisioning with \textbf{Terraform} as Infrastructure as Code (\textbf{30\%} faster deploys) and built a \textbf{Go + Terratest} framework catching \textbf{95\%} of misconfigurations before production.
\end{highlights}

\vspace{0.15 cm}

\begin{twocolentry}{Sep 2021 -- Dec 2021}
\textbf{Software Developer Intern} \,|\, Government of Maharashtra, Water Resource Department
\end{twocolentry}
\begin{highlights}
  \item Built and deployed a full-stack \textbf{React and MySQL} workflow application that digitized departmental record management and improved operational efficiency by \textbf{40\%}.
\end{highlights}

%═══════════════════════════════════════
\section{PROJECTS AND RESEARCH}
%═══════════════════════════════════════
\begin{twocolentry}{Jul 2026}
\textbf{NYC Mobility Forecasting} [PySpark, Apache Spark, Delta Lake, FastAPI, Redis]
\end{twocolentry}
\begin{highlights}
  \item Built Spark-based \textbf{medallion (bronze/silver/gold) pipeline} for large-scale batch data processing of a \textbf{2.3M-row} hourly dataset across \textbf{263 zones}; served predictions via a FastAPI service with PostgreSQL and Redis caching, containerized with Docker.
  \item Designed two-stage Poisson hurdle model improving weighted MAPE \textbf{18.77\% to 16.52\%} and cutting RMSE \textbf{43\%}, MAE \textbf{53\%} over baseline Random Forest.
\end{highlights}

\vspace{0.12 cm}

\begin{twocolentry}{Jul 2026}
\textbf{Spotify Analytics Platform} [Python, FastAPI, PostgreSQL, ETL, Materialized Views]
\end{twocolentry}
\begin{highlights}
  \item Engineered a 4-stage Python \textbf{ETL pipeline} for \textbf{71K+ streaming events} with rate limiting and backoff retry; designed a PostgreSQL schema with materialized views and 16 SQL RPC functions for low-latency aggregated queries.
\end{highlights}

%═══════════════════════════════════════
\section{EXTRACURRICULAR ACTIVITIES}
%═══════════════════════════════════════
\begin{onecolentry}
\begin{highlights}
  \item \textbf{Finalist, RevolutionUC 2025 Hackathon} (300 participants, 31 schools) \,|\, \textbf{SAFe Practitioner 5.1} (Mar 2023).
\end{highlights}
\end{onecolentry}

\end{document}
